\subsection*{S14. Array of scalars}

A field whose value is an array of scalars rather than of objects.

\begin{lstlisting}[style=json]
{ "id": 1, "tags": ["red", "blue"] }
\end{lstlisting}

Relationally this is S13 with a scalar element in place of an object: a child
table carrying a foreign key to the parent, an ordering index, and a single
value column. What is open is the module surface, since a record wrapping one
string is an awkward thing to hand a consumer.

\options
\begin{enumerate}[label=\textbf{(\alph*)}]
  \item \textbf{A full child module}, exactly symmetric with S13. Uniform, but
        the consumer unwraps a one-field record to reach a string.
  \item \textbf{Table without a module.} The accessor returns the scalars
        directly and no child module is generated. \selected
  \item \textbf{Both} --- generate the module, and add a convenience accessor
        returning the scalars.
  \item \textbf{Reject} arrays of scalars, requiring them to be wrapped as
        objects.
\end{enumerate}

\decision{The child table exists relationally, but generates no module. The
parent exposes a single accessor returning the scalars themselves.}

\begin{center}
\begin{tabular}{@{}l@{}}
\textbf{Table root} \\
\hline
\texttt{id} \\
\hline
1 \\
\hline
\end{tabular}
\qquad
\begin{tabular}{@{}llll@{}}
\multicolumn{4}{@{}l}{\textbf{Table tags}} \\
\hline
\texttt{id} & \texttt{root\_id} & \texttt{idx} & \texttt{value} \\
\hline
1 & 1 & 0 & red \\
2 & 1 & 1 & blue \\
\hline
\end{tabular}
\end{center}

\begin{lstlisting}[style=ml]
module Root : sig
  type id
  type t = { id : id }
  val get : db -> id -> t
  val all : db -> t list
  val tags : db -> t -> string list
end
\end{lstlisting}

\paragraph{Presence.} The presence marker of S13 applies unchanged. A partial
path yields \texttt{val tags : db -> t -> string list option}, with the parent
record carrying the presence marker of S21.3, and \texttt{Some []}
for a present empty array and \texttt{None} for an absent or null one.

\rationale{One construct rather than two. The consumer receives the
\texttt{string list} they actually want, with no one-field record to unwrap
and no second accessor to choose between, and the specification has a single
accessor per scalar-array path to describe rather than a pair whose
relationship must also be stated.

The relational layer is untouched: the table, its foreign key and its ordering
index all exist exactly as in S13. What is dropped is only the module that
would have wrapped them.}

\limitation{This breaks the design note's table-to-module correspondence: a
table exists with no module, so the correspondence holds in one direction
only. The rows of that table are also not addressable --- there is no
\texttt{Tags.id} and no way to name an individual element --- so if element
identity is ever needed, option (c) is the way back, and adopting it would be
additive rather than breaking.}
